% !TEX root = ../../main.tex

\section{Resultados}

Los resultados completos se obtienen desde los logs de resumen y de validación semántica generados para cada caso. En el estado actual del repositorio, los corpus sintéticos de \texttt{Maybe} y de \texttt{Dist.T Rational} están implementados y testeados; la tabla agregada final debe consolidar esos registros por familia y variante.

Como resultado representativo, el ejemplo académico probabilístico presentado a lo largo del informe reporta:

\begin{center}
\begin{tabular}{lccccc}
\hline
Caso & Sec. & ADo & Reord. & Perms. & Mínimas \\
\hline
Ejemplo académico & 4 & 3 & 2 & 5 & 4 \\
\hline
\end{tabular}
\end{center}

La tabla muestra la mejora estricta que motiva la memoria: el costo secuencial es 4, \texttt{ApplicativeDo} sobre el orden original reduce el costo a 3, y el reordenamiento conmutativo encuentra un plan de costo 2. Además, el caso genera cinco permutaciones válidas, de las cuales cuatro alcanzan el costo mínimo. El candidato 1, con orden \texttt{[1,3,2,4]}, es el primer mínimo y, por tanto, el seleccionado automáticamente.

La validación semántica ejecutó el programa original, \texttt{ApplicativeDo} sin reordenamiento, el candidato óptimo y cada una de las cinco permutaciones. Todas las variantes terminaron con código de salida cero y produjeron exactamente la misma distribución normalizada sobre preparación, quiz, dificultad y examen final. El registro \path{tests/main-example/semantic-validation.log} concluye con \texttt{RESULT: OK}.

Para el corpus sintético probabilístico, la tabla final debe reportar los mismos campos, junto con el resultado semántico de cada caso. La comparación debe realizarse sobre distribuciones normalizadas con \texttt{Dist.norm}. El control negativo con \texttt{IO} debe reportarse por separado, ya que su diferencia observable constituye el resultado esperado y no un fallo de los corpus evaluados bajo la hipótesis de conmutatividad.
